\section{TLA+ Specification of EventChain State Machine}\label{app:tla}

The TLA+ modules are maintained as source artifacts under \texttt{specs/}.
\texttt{specs/EventChain.tla} models a single ADL Lite concept as an append-only
sequence of events and formalises the well-formedness axioms from Definition~5
(\S\ref{subsec:compact-formal-spec}) together with the status lattice and confidence
G-Counter used in Theorems T1, T3, T4, T5, and T7.

\paragraph{State variables.}
\begin{itemize}
  \item \texttt{events}: a sequence of event records
    $(\text{event\_id}, \text{actor}, \text{event\_type}, \text{confidence}, \text{prev})$.
  \item \texttt{next\_id}: monotonic counter guaranteeing distinct event identifiers.
\end{itemize}

\paragraph{Actions.}
\begin{itemize}
  \item \texttt{AppendEvent(actor, etype, confidence)}: appends a new event whose
    \texttt{prev} field points to the previous event id, or $0$ for the genesis event.
\end{itemize}

\paragraph{Derived operators.}
\begin{itemize}
  \item \texttt{LatticeOrder} and \texttt{LUB}: the lifecycle partial order
    \texttt{provisional $<$ forked $<$ validated $<$ deprecated $<$ archived}.
  \item \texttt{DerivedStatus}: LUB of lifecycle event statuses in the chain.
  \item \texttt{DerivedConfidence}: maximum confidence among \texttt{VALIDATE} and
    \texttt{SNAPSHOT} events (G-Counter semantics).
\end{itemize}

\paragraph{Invariants checked by TLC.}
\begin{itemize}
  \item \texttt{Inv\_WellFormednessPreserved}: genesis anchoring, distinct ids,
    monotonic ids, non-empty actors, bounded confidence, valid event types.
  \item \texttt{Inv\_StatusMonotonic}: derived status is always a valid lattice
    element and confidence remains in $[0, \text{MaxConfidence}]$.
  \item \texttt{Inv\_MonotonicAppend}: every reachable status dominates
    \texttt{provisional} and confidence remains bounded.
\end{itemize}

\paragraph{Running TLC.}
A Python wrapper is provided at \texttt{scripts/run\_tlc.py}.  The default spec is
\texttt{EventChain}; the other supported specs are selected with \texttt{{-}-spec}:
\begin{verbatim}
python scripts/run_tlc.py --spec EventChain --max-events 20 --max-confidence 100
python scripts/run_tlc.py --spec CRDTMerge --max-events 10 --max-confidence 10
python scripts/run_tlc.py --spec ConsensusEngine --max-events 10 --max-confidence 10 --n-min 2
\end{verbatim}
The modules are intentionally bounded (\texttt{MaxEvents}, \texttt{MaxConfidence})
because TLC exhaustively explores a finite state space; the bound is a
model-checking limitation, not a system limitation. For \texttt{EventChain} with
$|\Sigma|=13$ event types, 2 actors, and $\text{MaxConfidence}=100$, TLC explores
$\approx 2.9 \times 10^7$ states at $\text{MaxEvents}=20$ (wall-clock time: 8.3~min on
Apple M2, 16~GB RAM). At $\text{MaxEvents}=25$, the state space exceeds $10^9$ and TLC
runs out of memory (16~GB), illustrating the expected exponential blow-up. The
bound of 20 events is therefore a practical model-checking limit; unbounded
correctness relies on the inductive argument given in Appendix~E and on the Coq
skeleton described in \S\ref{app:coq}.

\paragraph{CRDT merge spec.}
\texttt{specs/CRDTMerge.tla} models two concurrent branches of a single concept
that share a common genesis.  The \texttt{Merge} operator takes the set union of
both branches, applies last-writer-wins by \texttt{event\_id}, sorts by
\texttt{event\_id}, and re-anchors \texttt{prev} pointers.  TLC checks that the
merged chain is well-formed, that its status and confidence equal the LUB and
max of the branch values, and that merge is commutative, associative, and
idempotent (Theorem~9).

\paragraph{Consensus engine spec.}
\texttt{specs/ConsensusEngine.tla} models multiple agents appending to a shared
chain.  \texttt{Register} is allowed only on an empty chain; \texttt{Validate}
is guarded by distinct validators and an \texttt{N\_min} threshold; and
\texttt{Deprecate}/\texttt{Fork}/\texttt{Archive} are guarded by the ontology
lifecycle transition graph.  TLC checks well-formedness, valid transitions, the
minimum-validator requirement, and status/confidence boundedness.

\paragraph{Coq machine-verified proofs.}\label{app:coq}
A buildable Coq development is maintained under \texttt{formal/coq/}. As of the current revision, the following theorems have been completely machine-verified in Coq 8.18.0:

\begin{itemize}[leftmargin=2em]
  \item \textbf{T1 (Determinism of status derivation):} Proved in \texttt{Invariants.v} (205 lines). The proof shows that the status lattice is a join-semilattice and that the LUB over a totally ordered set is unique, hence $\delta(C)$ is uniquely determined for any well-formed chain.
  \item \textbf{E2 (Inductive status derivation correctness):} Proved in \texttt{Invariants.v}. The theorem \texttt{E2\_inductive\_status\_derivation} states that for any chain $C$ and event $e$, appending $e$ to $C$ computes the status incrementally as $\delta(C \oplus [e]) = \max(\delta(C), f(e.\tau))$, which equals the full LUB recomputation. Base cases for empty and singleton chains are also proved. This justifies the $O(1)$ incremental CRDT cache update in the Python implementation.
  \item \textbf{T3 (Status monotonicity):} Proved in \texttt{Status.v} (218 lines) and lifted to chains in \texttt{Invariants.v}. The proof establishes that the status lattice is a join-semilattice with the LUB operator, and that appending events preserves the monotonicity property: $\text{status\_leq}(\delta(C_1), \delta(C_1 \oplus C_2))$ for any chains $C_1, C_2$.
  \item \textbf{T4 (Confidence boundedness):} Proved in \texttt{Confidence.v} (153 lines) and lifted to chains in \texttt{Invariants.v}. The proof shows that $\gamma(C)$ is the max over a finite set of bounded values, hence $\gamma(C) \in [0, \text{MaxConfidence}]$.
  \item \textbf{T7 (Well-formedness preservation):} Proved in \texttt{Chain.v} (490 lines). The proof uses structural induction on the chain, with auxiliary lemmas for distinct-id preservation, monotonic-id preservation, and prev-linkage preservation under valid append. All 13 well-formedness axioms are formalised as Coq definitions; all six previously stubbed axioms (precondition evaluation, SHACL constraints, status transitions, lifecycle monotonicity, collusion resistance, and synthetic tagging) have been replaced with their full definitions and preservation proofs, introducing abstract Ed25519/SHA-256 cryptographic primitives in \texttt{Crypto.v} (77 lines). The Coq development now contains zero remaining \texttt{Admitted} lemmas.
  \item \textbf{T9 (CRDT merge convergence):} Proved in \texttt{CRDT.v} (885 lines). The file formalises the event-set merge operation, last-writer-wins resolution, and re-anchoring of prev pointers. The commutativity, associativity, and idempotency properties are proved, as well as well-formedness preservation under merge (propagating scope ACL, signature verification, and confidence clamping through the re-anchor operation).
\end{itemize}

The Coq development totals $\sim$2,500 lines across seven modules. All algebraic proofs are closed with zero remaining \texttt{Admitted} lemmas. The TLA\textsuperscript{+} specification (\texttt{specs/}) has been syntax-checked and model-checked for chains up to 20 events, complementing the Coq deductive proofs with finite-state verification.

Optional Iris stubs in \texttt{iris/} provide a placeholder resource algebra and a Hoare-triple skeleton for the split-lock \texttt{append} operation; this is deferred to future work (FW3).
